\subsection*{A polynomial cutoff for the literal source}
The small physical endpoint need not force an exponentially large Fourier
cutoff for this fixed two-mode source.  The additional ingredient is an
exterior norm estimate for a finite entire extension of the arithmetic
sum.  It controls the endpoint derivative jets without bounding the whole
interior variation by the endpoint amplitude.

\begin{lemma}[Polynomial approximation with a growing derivative order]
\label{lem:v114-entire-polynomial}
Fix a compact real interval $I$ of positive length. Suppose
\[
 F(z)=\int_{-\Omega}^{\Omega}e^{izt}\,d\mu(t),\qquad
 \Omega\le A_0c,\quad \norm\mu\le c^{A_1}.
\]
For each fixed $H>0$, there is $A>0$ such that a polynomial $T$ of degree
$D=\lceil Ac\rceil$ satisfies, simultaneously for $0\le j\le\lceil c\rceil$,
\begin{equation}\label{eq:v114-polynomial-error}
 \norm{(F-T)^{(j)}}_{L^\infty(I)}
       \le e^{-Hc}(Cc)^j.
\end{equation}
All constants are independent of $c$ and $j$. Moreover,
\[
 \norm T_\infty\le C(D+1)\norm T_2,\qquad
 \norm{T^{(j)}}_\infty\le(CD^2)^j\norm T_\infty
\]
on $I$.
\end{lemma}
\begin{proof}
Take the degree-$D$ Taylor polynomial about the midpoint of $I$ under the
finite measure. If $d$ is the half-length of $I$, its differentiated
remainder is at most
\[
 \norm\mu\,\Omega^j e^{\Omega d}
       \frac{(\Omega d)^{D+1-j}}{(D+1-j)!}.
\]
The inequality $k!\ge(k/e)^k$, with $A$ sufficiently large, proves
\eqref{eq:v114-polynomial-error} for the entire stated range of $j$.
The first polynomial inequality follows from normalized Legendre expansion,
$|P_k|\le1$, and $\sum_{k=0}^D(2k+1)=(D+1)^2$. The second is Markov's
inequality iterated on the derivatives, with the fixed interval scaling.
Real and imaginary parts give the same bounds for complex polynomials.
\end{proof}

\begin{lemma}[Endpoint jets of the finite arithmetic sum]
\label{lem:v114-endpoint-jets}
Put $c=2\pi\lambda^2$, $r=\lceil c\rceil$, and let $p_\lambda^\circ$
be the inversion-even proxy before the fixed value-at-zero repair.
If $J_j$ is the sum of absolute jumps of its $j$th ordinary logarithmic
derivative, including the jump at the identified periodic endpoint, then
\begin{equation}\label{eq:v114-jet-bound}
 J_j\le C|B_\lambda^+|\lambda^3(C\lambda^4)^j,
          \qquad 1\le j<r.
\end{equation}
Its zeroth-order periodic boundary jump is zero. On its smooth pieces,
for fixed constants $A,C,\rho>0$,
\begin{equation}\label{eq:v114-smooth-remainder}
 \norm{(p_\lambda^\circ)^{(r)}_{\rm ac}}_1
       \le CL\lambda^3 r!\rho^{-r}e^{Ac}.
\end{equation}
\end{lemma}
\begin{proof}
Write
\[
 g(x)=\sqrt\lambda h_\lambda^\circ(\lambda x)
     =\sum_{j=0,4}a_{j,\lambda}\psi_{j,c}(x),\qquad g(1)=B_\lambda^+.
\]
The exact finite Fourier eigenrelation used in
Lemma~\ref{lem:v12-endpoint-energy} represents $g$ by a measure supported
on $[-c,c]$ with total variation $O(\lambda)$; this is also the
bandlimited extension in \cite[\S30.15]{DLMFPSWF}.
Proposition~\ref{prop:endpoint-radial} and endpoint noncancellation give
$|g(x)|\le C|B_\lambda^+|$ for $x\ge1$.

Set $M=\lceil2\lambda^2\rceil$ and
$G(v)=\sum_{m=1}^M g(mv/\lambda^2)$. Its measure has bandwidth
$2\pi M=O(c)$ and total variation $O(\lambda^3)$. For $1/2\le v\le1$,
\[
 G(v)=\lambda v^{-1/2}f_\lambda^\circ(v/\lambda)
       +\sum_{\substack{m\le M\\mv\ge\lambda^2}}
                       g(mv/\lambda^2)
\]
almost everywhere. The last sum is $O(|B_\lambda^+|\lambda^2)$.
Corollary~\ref{cor:pswf-cross-tail}, which includes the unrepaired source,
and $dv/v=du/u$ therefore give
\[
 \norm G_{L^2(1/2,1)}\le C|B_\lambda^+|\lambda^2.
\]
Since $|B_\lambda^+|\asymp c^{11/4}e^{-c}$, apply
Lemma~\ref{lem:v114-entire-polynomial} with fixed sufficiently large $H$.
Its polynomial inequalities imply
\begin{align*}
 |G^{(j)}(1)|&\le C|B_\lambda^+|\lambda^4(C\lambda^4)^j,\\
 \sup_{1\le x\le3}|g^{(j)}(x)|
       &\le C|B_\lambda^+|(C\lambda^4)^j,
          \qquad 0\le j\le r.
\end{align*}
For the second bound approximate $g$ on $[1,3]$ and use its already
bounded sup norm, so the extra Legendre degree factor is unnecessary.

At the strict interior lower trace $v=1+$, subtract from $G$ precisely the
terms with $m\ge\lambda^2$, including equality when $\lambda^2$ is an
integer, and multiply by $\sqrt v/\lambda$. Their arguments lie in
$[1,3]$, their number is $O(\lambda^2)$, and differentiation costs at most
$3^j$. Thus the logarithmic trace jets are bounded by
$C|B_\lambda^+|\lambda^3(C\lambda^4)^j$.
To justify the conversion uniformly in $j$, use
$(v\partial_v)^j=\sum_k S(j,k)v^k\partial_v^k$ and
$S(j,k)\le\binom jk j^{j-k}$; hence
$\sum_k S(j,k)A^k\le(A+j)^j$. Here $j=O(\lambda^2)$ and
$A=C\lambda^4$. For the product factor, $|(\sqrt v)^{(k)}|_{v=1}\le k!$ and
$\sum_{k=0}^j\binom jk k!A^{j-k}\le2A^j$ when $A\ge2j$;
thus its cost is uniform in $j$ as well.
At the upper trace only the first summand survives, giving the stronger
bound without $\lambda^3$.

Each interior derivative jump is a source endpoint jet times
$m^{-1/2}$, with the reflected sign. Summing $m<\lambda^2$ costs
$O(\lambda)$; the preceding trace estimates control the periodic jump.
This proves \eqref{eq:v114-jet-bound}. Inversion symmetry gives the zero
function jump at the periodic boundary.
Finally extend each smooth piece using its fixed finite collection of
entire summands. On a complex logarithmic disk of fixed small radius
$\rho$, every active source argument has imaginary part bounded by a
fixed constant. The Fourier representation bounds the extension by
$C\lambda^3e^{Ac}$. Cauchy's estimate, integrated over total length $L$,
proves \eqref{eq:v114-smooth-remainder}. No asymptotic remainder was
differentiated, and the smooth compact repair is not used in this step.
\end{proof}

\begin{proposition}[Polynomial physical-endpoint recovery]
\label{prop:v114-polynomial-endpoint}
For the exact repaired source in Proposition~\ref{prop:v12-source}, set
\[
 N_\lambda=\left\lceil\lambda^8(1+L)\right\rceil,
       \qquad k_\lambda=P_{N_\lambda}p_\lambda.
\]
For all sufficiently large $\lambda$,
\begin{align}
 \frac{|\delta_{N_\lambda}(k_\lambda)-B_\lambda|}{|B_\lambda|}
       &\le\frac C\lambda,\label{eq:v114-polynomial-endpoint}\\
 |\widehat p_\lambda(n)|
       &\le C|B_\lambda|\frac{\lambda\sqrt L}{|n|},
            \qquad |n|>N_\lambda.\label{eq:v114-high-tail}
\end{align}
This is an asymptotic estimate with no certified numerical threshold or
optimized exponent. The vector is the literal ordinary Fourier projection;
no endpoint shell is added.
\end{proposition}
\begin{proof}
First omit the compact repair. The periodic proxy is BV and continuous
at the identified endpoint, so the Dirichlet--Jordan theorem identifies
its symmetric Fourier sum there with its common one-sided trace. We may
therefore estimate the actual endpoint error by its symmetric Fourier tail.
In the jump train of
Lemma~\ref{lem:g1-jump-train}, the endpoint tail contributed by a threshold
$m$ is, up to sign,
\[
 \frac{B_\lambda^+}{\pi\sqrt m}
       \sum_{n>N}\frac{\sin(2\pi n\log m/L)}n.
\]
Dirichlet summation and the bounded sine-series partial sums give
\[
 \left|\sum_{n>N}\frac{\sin(n\theta)}n\right|
       \le C\min\left(1,\frac1{N|\sin(\theta/2)|}\right).
\]
Separate the largest integer $m_0<\lambda^2$: its contribution is
$O(|B_\lambda^+|/\lambda)$ even when $\lambda^2$ approaches an integer.
For every remaining $m$, $\lambda^2-m\ge1$. Splitting at $m=\lambda$
and $m=\lambda^2/2$ gives
\[
 \sum_{\substack{1<m<\lambda^2\\m\ne m_0}}
 \frac{m^{-1/2}}{|\sin(\pi\log m/L)|}
       \le C\lambda L(1+\log\lambda).
\]
Near the upper end use
$\log(\lambda^2/m)\ge(\lambda^2-m)/\lambda^2$ and sum the harmonic
denominator. Thus the entire function-jump endpoint error divided by
$|B_\lambda^+|$ is at most
\begin{equation}\label{eq:v114-resonance-bound}
 C\left\{\lambda^{-1}+\frac{\lambda L(1+\log\lambda)}N\right\}.
\end{equation}
This retains, rather than assumes away, near-threshold arithmetic resonance.

Integrate by parts $r=\lceil c\rceil$ times on the smooth pieces, retaining
the jumps of every derivative and the periodic derivative jumps. With
$t_n=2\pi n/L$, the $j$th jump term in a normalized Fourier coefficient
has magnitude at most $J_j/(\sqrt L\,|t_n|^{j+1})$. Its physical-endpoint
tail is at most
\[
 \frac{2J_j}{L}\sum_{n>N}|t_n|^{-j-1}
       \le \frac{C J_j}{j}\left(\frac L{2\pi N}\right)^j.
\]
By Lemma~\ref{lem:v114-endpoint-jets}, summing $1\le j<r$ gives
$C|B_\lambda^+|\lambda^7L/N$ whenever $C\lambda^4L/N<1/2$.
The final smooth-piece remainder, after endpoint summation and division by
$|B_\lambda^+|$, is bounded by
\begin{equation}\label{eq:v114-analytic-tail}
 C\lambda^A(1+L)^Ae^{Ac}r!
                 \left(\frac{CL}{N}\right)^{r-1}.
\end{equation}
Here $A,C$ are fixed and independent of $r,\lambda$.
At $N=N_\lambda$, $r!\le(C\lambda^2)^r$, so the logarithm of this bound
is at most $-6c\log\lambda+O(c+\log\lambda+\log L)$.
It is smaller than every fixed inverse power of $\lambda$.
Combining this with \eqref{eq:v114-resonance-bound} proves the first
claim for $p_\lambda^\circ$.

The same expansion gives, at every $|n|>N_\lambda$,
\[
 |\widehat p_\lambda^\circ(n)|
 \le \frac{C|B_\lambda^+|}{\sqrt L}
       \left(\frac\lambda{|t_n|}+\frac{\lambda^7}{|t_n|^2}\right)
       +|R_n|.
\]
The second term and remainder are absorbed by the first. For the remainder,
its ratio to $|B_\lambda^+|\lambda/(\sqrt L|t_n|)$ decreases with $|n|$
and is already smaller than every fixed inverse power at $N_\lambda$.
This proves \eqref{eq:v114-high-tail} before repair.

Finally let $e=p_\lambda-p_\lambda^\circ$ and
$a=h_\lambda^\circ(0)$. Direct counting gives
$\norm e_\infty\le C\lambda^{3/2}|a|$.
The Fourier Lebesgue constant is $O(\log(N+2))$, so
\[
 \frac{|(P_N-I)e(0)|}{|B_\lambda|}
 \le C\lambda^{3/2}\log(N+2)\frac{|a|}{|B_\lambda|}
       =O_A(\lambda^{-A})
\]
along this polynomial diagonal, by \eqref{eq:v12-zero-value}.
For the coefficient assertion a direct variation estimate is stronger:
if the fixed repair bump is supported in $[-b,b]$, then
\[
 \int_{1/\lambda}^{\lambda}|D_{\log}\mathcal E(\psi)(u)|\frac{du}{u}
 \le \sum_{m\le b\lambda}m^{-1/2}
   \int_0^b x^{-1/2}\bigl(\tfrac12|\psi(x)|+x|\psi'(x)|\bigr)\,dx
 \le C\sqrt\lambda.
\]
Inversion averaging preserves this bound and has no periodic function
jump. Hence $|\widehat e(n)|\le C|a|\sqrt{\lambda L}/|n|$, which is
negligible. Use $B_\lambda^+\asymp B_\lambda$ to finish.
\end{proof}

\begin{corollary}[Weak G1 on the same polynomial-cutoff vector]
\label{cor:v114-polynomial-g1}
Let $R>3$, $N=N_\lambda$, $K\ge N$, and $x\in E_K$, measured by
the weighted norm of Proposition~\ref{prop:g1-weighted-sampler}.
For all sufficiently large $\lambda$,
\begin{align}
 \frac{|\ip{D_{\log}W_{\lambda,K}k_\lambda}{x}|}
      {|\delta_K(k_\lambda)|}
 &\le C_R\norm x_{\mathcal H_{R,L}}
 \left[\mathfrak r(\lambda)+\lambda^2
 \left\{\frac LN+\left(\frac LN\right)^{R-1}
                   (1+\log(2N))\right\}\right].
 \label{eq:v114-polynomial-g1}
\end{align}
The bound is uniform in $K$ and tends to zero on bounded weighted test
families. This is a fixed-source weak pairing statement, not control
in operator norm on a growing source block.
\end{corollary}
\begin{proof}
The coefficients of $(P_N-I)p_\lambda$ vanish for $|n|\le N$.
Insert \eqref{eq:v114-high-tail} into the absolute convolution proof
of Proposition~\ref{prop:g1-weighted-sampler}: on the only frequencies
used there, $\mathcal V_\lambda$ is replaced by $C|B_\lambda|\lambda$.
Both its off-diagonal estimate and its separately treated logarithmic
diagonal therefore give the displayed $\lambda^2$ error. The sums
converge absolutely for $R>3$; the fixed-window domain identification
in Proposition~\ref{prop:v114-weil-core} also justifies the passage
from finite Fourier sums to $p_\lambda$.
The continuum term is Proposition~\ref{prop:g1-sobolev}. Finally
$|\delta_K(k_\lambda)|\ge(1-C/\lambda)|B_\lambda|$ and, with the
first-slot-linear convention,
$\ip{D_{\log}W_{\lambda,K}k_\lambda}{x}
=QW_\lambda(k_\lambda,D_{\log}x)$.
Unlike the filtered vector, $P_Np_\lambda$ lies in $E_N$, so $K\ge N$
suffices.
\end{proof}

These results assert no positive Weil complement along this diagonal,
a small inverse-weighted spectral residual, or control of the corrected
sampler's graph defect. Those obligations remain separate. The fixed
orders $0,4$ are essential; no uniform claim over a growing prolate block
is made.
